\chapter{Bibliography and Reading Order}
\label{ch:app_biblio}
\chaptermark{Bibliography and reading order}

This appendix reorganises the bibliography as three suggested reading
orders, depending on your goal. If you have read cover to cover, you
have used the default; the orderings below are for readers who want to
focus. It also serves as an annotated bibliography: every source cited
anywhere in the book is collected here, grouped by topic, with a
one-sentence note on why you would read it and which chapters draw on
it. The final two sections list the primary sources that are free
online and the topics this book deliberately does not cover.

\begin{backgroundbox}[How to use this appendix]
There are three kinds of reader this book is written for: someone
preparing for the IMC Prosperity competition, someone starting an
internship or graduate role on a bank Strats desk (Goldman Sachs or
otherwise), and a completionist who wants the full picture. The
\emph{Three reading orders} below (\cref{sec:app_biblio_tracks}) give
the first two a fast route through the material that matters for their
goal. The \emph{Annotated bibliography}
(\cref{sec:app_biblio_annotated}) is the reference side: fifteen
topic clusters, each listing the specific books, papers, and blog
posts that feed the relevant chapters, with a short note on why the
source is worth reading. The \emph{Free online sources}
(\cref{sec:app_biblio_free}) section lists the material you can legally
read tonight without paying for anything. Finally,
\emph{What's not in this book} (\cref{sec:app_biblio_gaps}) is an
honest list of topics this book avoids, with pointers for where to go
next.
\end{backgroundbox}

% =========================================================================
\section{Three reading orders}
\label{sec:app_biblio_tracks}
% =========================================================================

Reading 635 pages cover to cover is a real commitment. If you picked
this book up with a specific goal in mind, one of the three tracks
below will get you there faster. Each track is a sequence of chapters
with the pages that would otherwise be in the way skipped. The tracks
overlap substantially in the foundations; they diverge in what they
emphasise after \cref{ch:options_payoffs}.

\subsection{Track 1 --- Prosperity competition}

The Prosperity track is about producing a working algorithm in five
rounds. You need enough microstructure to understand why the simulator
behaves the way it does, the small handful of strategies that keep
winning writeups (market making, pairs/basket arbitrage, options
scalping, voucher stat arb), and a playbook for diagnosing which
strategy to apply to each product. You do \emph{not} need rates,
credit, commodities, SecDB, FVA, or regulatory capital.

\begin{enumerate}
  \item \textbf{Orientation} (Ch 0--3). How to read the book, what a
        market is, price and value, the mechanics of an order book and
        matching engine. Skipping this is the single most common
        reason Prosperity writeups are full of confused terminology.
  \item \textbf{Microstructure minimum} (Ch 4--7). Bid--ask spreads,
        the Roll/Kyle/Glosten--Milgrom models, inventory dynamics, and
        market impact at an intuitive level. You do not need the full
        Hasbrouck machinery for Prosperity --- a few hundred ticks is
        not enough data for it anyway --- but the vocabulary matters.
  \item \textbf{Core strategies} (Ch 8--13). Market making and
        Avellaneda--Stoikov, mean reversion and Ornstein--Uhlenbeck,
        pairs trading and cointegration, basket/ETF stat arb,
        momentum, arbitrage zoo. These six chapters are the
        \emph{meat} of what wins Prosperity rounds.
  \item \textbf{Options minimum} (Ch 14--17). Just enough to trade the
        volcanic rock vouchers: option payoffs, Black--Scholes, the
        Greeks, and delta hedging. You can skip the full vol surface
        chapter on a first pass and come back to it for Round 5.
  \item \textbf{Which-strategy-when} (Ch 22). The pattern-recognition
        chapter is the keystone for Prosperity: given a new product on
        day one of a round, how do you decide what to do with it?
  \item \textbf{Prosperity playbook} (Ch 23--26). Round-by-round
        analysis, the Frankfurt Hedgehogs' 2nd-place writeup as a case
        study, trader-ID exploitation, and manual trading tricks. This
        is where all the theory cashes out into leaderboard points.
\end{enumerate}

Roughly 250 pages. Prerequisite maths: comfortable with probability at
the level of expected value, variance, and the normal distribution;
first-year calculus; and enough linear algebra to invert a $2\times 2$
matrix. No stochastic calculus required.

\subsection{Track 2 --- Goldman Sachs Strats}

The GS track is about walking onto a trading floor and not being
useless in week one. That means knowing the business (what equities,
rates, credit, commodities, FX desks actually trade), the quant
machinery (Black--Scholes rigorously, vol surface, Greeks, numerical
methods), execution and risk in production (limits, VaR, liquidity
risk), and the institutional vocabulary (SecDB, Slang, FVA, Volcker,
MiFID). You \emph{do} want the interview and culture material; you do
not need the full Prosperity playbook.

\begin{enumerate}
  \item \textbf{Foundations} (Ch 1--3). Markets, price formation,
        order books. Same as Track 1.
  \item \textbf{Microstructure} (Ch 5--7). The Kyle, Glosten--Milgrom
        and market-impact chapters. On a rates or credit desk you
        will not be quoting in a limit order book, but the
        information-asymmetry framework is part of how every quant
        thinks about markets.
  \item \textbf{Options and volatility} (Ch 14--17). Black--Scholes
        derived from first principles, the Greeks, delta and gamma
        hedging, and the full volatility surface (skew, term
        structure, sticky-delta vs sticky-strike). This is the
        \emph{minimum} a bank Strat must know.
  \item \textbf{Execution, risk, and methods} (Ch 18--22). Execution
        algorithms (VWAP, Almgren--Chriss), slippage, risk management
        (VaR, stress, Kelly), backtesting hygiene, ML for trading,
        and the pattern-recognition chapter.
  \item \textbf{Asset-class depth} (Ch 27--33). Rates and fixed
        income, credit derivatives, FX, commodities, structured
        products, and a broad survey of how desks are organised.
  \item \textbf{GS-specific} (Ch 34--38). Numerical methods in the
        Strats toolkit, production risk systems, SecDB and Slang,
        regulatory capital and Volcker, and the interview-problem
        chapter.
\end{enumerate}

Roughly 400 pages. Prerequisite maths: the whole of \cref{ch:app_math},
including stochastic calculus to the level of It\^o's lemma and SDEs.
If your probability is rusty, spend a weekend with Grimmett and
Stirzaker before Chapter 14.

\subsection{Track 3 --- Full cover-to-cover}

The default. Read in order. 635 pages plus appendices. The only
reason to pick a track over this one is time. Everything in Tracks 1
and 2 is in here; the extra chapters you pick up on the full path are
the ones you will be glad of the \emph{second} time you need a
reference.

\begin{keyfactbox}[Three reading tracks at a glance]
\begin{center}
\small
\begin{tabular}{@{}l l c l@{}}
\toprule
\textbf{Track} & \textbf{Chapters} & \textbf{Pages} & \textbf{Maths level} \\
\midrule
Prosperity         & 0--17 (core), 22, 23--26       & $\sim\!250$ & Prob + calc  \\
GS Strats          & 1--3, 5--7, 14--22, 27--38     & $\sim\!400$ & Prob + stoch calc \\
Full cover-to-cover & 0--38 + App A--C              & $\sim\!635$ & All of App A \\
\bottomrule
\end{tabular}
\end{center}

\noindent\emph{Prob + calc} means you are comfortable with
expectation, variance, normal distributions, and single-variable
calculus. \emph{Prob + stoch calc} adds It\^o's lemma, stochastic
differential equations, and risk-neutral pricing at the level of
Shreve's book. \emph{All of App A} means all of the above plus the
PDE and numerical linear algebra material.
\end{keyfactbox}

% =========================================================================
\section{Annotated bibliography by topic}
\label{sec:app_biblio_annotated}
% =========================================================================

This section reorganises the references by topic rather than
alphabetically. Every source given a \verb|\cite| macro in
\texttt{conventions.tex} is covered, and every bibliographic entry in
\texttt{references.bib} is cited at least once in one of the topic
clusters below. Each entry lists the source, a one-sentence note on
why you would read it, and pointers to the chapters that cite it.

\subsection{Market microstructure}
\label{sec:biblio_microstructure}

\begin{description}
  \item[\citet{bouchaud2018trades} --- \emph{Trades, Quotes and Prices}.]
    The modern book on limit order book dynamics, written by the
    Capital Fund Management team; empirically grounded, with chapters
    on square-root market impact, order flow autocorrelation, and
    Hawkes processes. Cited in Chs~4--7, 12, 18.
  \item[\citet{harris2003trading} --- \emph{Trading and Exchanges}.]
    The single best ``how markets actually work'' textbook, aimed at
    practitioners; no equations, exhaustive on institutional detail
    (makers, takers, fees, tick sizes, exchange rulebooks). Cited in
    Chs~1--4, 5, 7.
  \item[\citet{hasbrouck2007empirical} --- \emph{Empirical Market
    Microstructure}.] The econometrician's companion to Harris;
    teaches you how to estimate spread decompositions, VAR models of
    trade and quote dynamics, and Hasbrouck information shares. Cited
    in Chs~4, 5, 7.
  \item[\citet{foucault2024market} --- \emph{Market Liquidity}.]
    Foucault, Pagano and R\"oell, 2nd edition. Theoretical treatment
    of liquidity and market design, covering fragmentation, tick-size
    reforms, and dark pools. Cited in Chs~5, 7, 18.
  \item[\citet{roll1984simple} --- Roll's implicit spread estimator.]
    A beautifully simple paper: estimate the effective bid--ask spread
    from the serial covariance of returns. Cited in \cref{ch:spreads_roll}
    (Ch~4).
  \item[\citet{kyle1985continuous} --- Kyle's continuous auction
    model.] The foundational paper on strategic informed trading and
    price impact; where ``Kyle's lambda'' comes from. Cited in Chs~5,
    18.
  \item[\citet{glosten1985bid} --- Glosten and Milgrom.] The
    information-asymmetry dealer model: why spreads exist and how
    market makers protect themselves against adverse selection.
    Cited in \cref{ch:kyle_gm} (Ch~5).
  \item[\citet{hawkes1971spectra} --- Hawkes processes.] Self-exciting
    point processes; the modern tool for modelling order-arrival
    clustering. Cited in Chs~7, 18.
  \item[\citet{pelgrims2020matching} --- Pelgrims, \emph{Matching
    Engines}.] A clear, engineering-focused blog post on how a
    price-time priority matching engine is actually implemented.
    Cited in \cref{ch:lob} (Ch~3).
\end{description}

\subsection{Market making}

\begin{description}
  \item[\citet{avellaneda2008high} --- Avellaneda and Stoikov.] The
    paper behind the market-making quoting formula that every modern
    HFT shop's textbook starts from; inventory-aware bid/ask skew
    under exponential arrival intensity. Cited in Chs~8, 22.
  \item[\citet{hostoll1981} --- Ho and Stoll.] The original
    inventory-based dealer pricing model; A--S is the continuous-time
    refinement of this. Cited in \cref{ch:market_making}.
  \item[\citet{hummingbotASguide} --- Hummingbot's A--S guide.]
    A practitioner's walk-through of Avellaneda--Stoikov with the
    parameter intuition laid out in plain English; the best first
    read before tackling the original paper. Cited in Chs~8, 23.
\end{description}

\subsection{Mean reversion and pairs trading}

\begin{description}
  \item[\citet{chan2013algorithmic} --- Chan, \emph{Algorithmic
    Trading}.] A compact, practitioner-focused book with working
    Python for mean-reversion, momentum, and stat-arb strategies;
    the single best bridge between academic theory and a deployable
    algo. Cited in Chs~9--11, 19--22.
  \item[\citet{chan2022quantitative} --- Chan, \emph{Quantitative
    Trading}.] The earlier, more business-facing Chan book; more
    on infrastructure, less on strategy details. Cited in Chs~11, 19.
  \item[\citet{cartea2015algorithmic} --- Cartea, Jaimungal and
    Penalva.] The rigorous textbook for execution and market making
    under stochastic control; includes the full Avellaneda--Stoikov
    derivation and extensions. Cited in Chs~8, 9, 18.
\end{description}

\subsection{Basket and statistical arbitrage}

\begin{description}
  \item[\citet{avellaneda2010stat} --- Avellaneda and Lee.] The
    canonical PCA-residual stat-arb paper; how to isolate
    idiosyncratic mean-reverting signals from a large cross-section.
    Cited in \cref{ch:basket_etf_arb} (Ch~11).
  \item[\citet{frankfurthedgehogs2025} --- Frankfurt Hedgehogs,
    Prosperity 3.] The 2nd-place 2025 Prosperity writeup; the
    clearest public description of how to win the basket and
    voucher arbitrage rounds. Cited in Chs~11, 23--26.
  \item[\citet{frankfurtHedgehogs2025imcP3} --- Frankfurt Hedgehogs
    polished algorithm.] The companion code repository to the
    writeup above. Cited in Chs~23--26.
\end{description}

\subsection{Momentum and trend following}

\begin{description}
  \item[\citet{moskowitz2012time} --- Moskowitz, Ooi and Pedersen.]
    The ``time-series momentum'' paper --- the empirical backbone of
    modern CTA strategies. Cited in \cref{ch:momentum} (Ch~12).
  \item[\citet{jegadeesh1993} --- Jegadeesh and Titman.] The classic
    cross-sectional momentum paper; together with MOP it forms the
    evidence base for systematic trend strategies. Cited in Ch~12.
  \item[\citet{daniel2016momentum} --- Daniel and Moskowitz,
    \emph{Momentum Crashes}.] The complement to MOP: documents when
    momentum fails catastrophically and why. Cited in Chs~12, 19.
\end{description}

\subsection{Arbitrage}

\begin{description}
  \item[\citet{damodaranOptionsArbitrage} --- Damodaran's options
    arbitrage notes.] A free NYU Stern note that collects put--call
    parity, conversion/reversal, and box spreads in one accessible
    document. Cited in Chs~13, 15.
  \item[\citet{stoll1969relationship} --- Stoll's put--call parity
    paper.] The original derivation of put--call parity. Cited in
    \cref{ch:arbitrage_zoo} (Ch~13).
  \item[\citet{merton1973rational} --- Merton's rational option
    pricing.] Extends Black--Scholes to a proper no-arbitrage
    framework; contains the American-option arbitrage bounds. Cited
    in Chs~13, 14.
\end{description}

\subsection{Options and Black--Scholes}

\begin{description}
  \item[\citet{hull8ed} --- Hull, \emph{Options, Futures, and Other
    Derivatives}.] The industry-standard derivatives textbook;
    encyclopaedic, readable, and the reference for anything you need
    to revise before a GS interview. Cited throughout Chs~14--17,
    27--32.
  \item[\citet{natenberg2015option} --- Natenberg, \emph{Option
    Volatility and Pricing}.] The options-trader's bible; everything
    about how vol actually trades (skew, term structure, events),
    with almost no equations. Cited in Chs~14--17.
  \item[\citet{mit18S096} --- MIT 18.S096, lectures 15--21.] Free MIT
    OCW lectures that rigorously derive It\^o calculus and
    Black--Scholes from first principles; the best free option if
    Shreve feels too heavy. Cited in Chs~14, 16, 34.
  \item[\citet{black1976pricing} --- Black's 1976 commodity-options
    paper.] Introduces the Black-76 variant used for futures options
    and interest-rate derivatives. Cited in Chs~28, 30.
  \item[\citet{coxrossrubinstein1979} --- CRR binomial trees.]
    The discrete-time binomial pricing model; the pedagogical
    counterpart to Black--Scholes. Cited in Chs~14, 34.
\end{description}

\subsection{Volatility surface, gamma, and vol trading}

\begin{description}
  \item[\citet{bennett2014trading} --- Bennett, \emph{Trading
    Volatility}.] A 250-page free PDF from a sell-side vol desk;
    correlation, term structure, skew, dispersion, variance swaps.
    Absurdly practical. Cited in Chs~16, 17.
  \item[\citet{sinclair2013volatility} --- Sinclair, \emph{Volatility
    Trading}.] The retail-trader's companion to Bennett; smaller,
    more opinionated, excellent on dispersion and vol arbitrage.
    Cited in Ch~17.
  \item[\citet{dupire1994pricing} --- Dupire's local vol.] Derives
    the unique local volatility surface consistent with observed
    option prices. Cited in Ch~17.
  \item[\citet{derman1994riding} --- Derman and Kani.] The parallel
    development of local vol; the ``sticky-strike vs sticky-delta''
    intuition originates here. Cited in Ch~17.
  \item[\citet{heston1993closed} --- Heston stochastic volatility.]
    The closed-form stoch-vol model; the baseline for every more
    complex vol model used on a desk. Cited in Chs~17, 34.
  \item[\citet{demeterfi1999guide} --- Demeterfi et al., GS QS
    research note on variance swaps.] The paper that made variance
    swaps standard; replicate variance with a weighted strip of
    options. Cited in Chs~17, 34.
\end{description}

\subsection{Execution}

\begin{description}
  \item[\citet{almgren2000optimal} --- Almgren and Chriss.] The
    foundational paper on optimal execution under a linear impact
    model; where VWAP-style schedules come from. Cited in Chs~18,
    22.
  \item[\citet{almgren2005direct} --- Almgren, Thum, Hauptmann, Li.]
    Empirical calibration of the square-root market-impact law on
    Citigroup equity data; the paper you cite when a PM asks ``how
    big can I trade?''. Cited in Ch~18.
\end{description}

\subsection{Risk management and Kelly}

\begin{description}
  \item[\citet{thorp2006kelly} --- Thorp's Kelly survey.] Ed Thorp's
    own summary of Kelly sizing across blackjack, sports betting, and
    the stock market; the clearest introduction to fractional-Kelly
    practice. Cited in Chs~19, 22.
  \item[\citet{kelly1956new} --- Kelly's original paper.] The 1956
    Bell Labs paper that started it all; short, derives $f^* = p -
    q/b$ from information-theoretic first principles. Cited in Ch~19.
  \item[\citet{lowenstein2000ltcm} --- Lowenstein, \emph{When Genius
    Failed}.] The LTCM narrative history; essential reading for why
    leverage and correlation can kill even Nobel laureates. Cited in
    Chs~19, 35.
\end{description}

\subsection{Backtesting and machine learning}

\begin{description}
  \item[\citet{lopezdeprado2018afml} --- L\'opez de Prado,
    \emph{Advances in Financial Machine Learning}.] The rigorous
    book on why naive ML for finance fails (leakage, overfitting,
    non-IID data) and what to do about it (triple-barrier labels,
    purged CV, fractional differentiation). Cited in Chs~20, 21.
  \item[\citet{bailey2014deflated} --- Bailey and L\'opez de Prado,
    deflated Sharpe.] Shows how Sharpe ratios lie when you select
    from many trials; essential reading before publishing any
    backtest. Cited in Chs~20, 21.
  \item[\citet{jansen2020ml} --- Jansen, \emph{Machine Learning for
    Algorithmic Trading}.] The 900-page how-to; more code, less
    rigour than L\'opez de Prado; great for feature engineering
    recipes. Cited in Ch~21.
  \item[\citet{chan2020cpo} --- Chan's conditional parameter
    optimisation blog.] A short note on adapting strategy parameters
    to market regimes. Cited in Chs~21, 22.
  \item[\citet{quantstart2016hmm} --- QuantStart HMM regime detection.]
    A hands-on blog post on fitting Hidden Markov Models to equity
    returns. Cited in Ch~22.
  \item[\citet{rabiner1989tutorial} --- Rabiner's HMM tutorial.]
    The standard reference for HMMs; forward--backward, Viterbi,
    Baum--Welch. Cited in Ch~22.
\end{description}

\subsection{Prosperity writeups and tooling}

\begin{description}
  \item[\citet{frankfurthedgehogs2025} --- Frankfurt Hedgehogs.]
    2nd-place 2025 writeup (see \cref{sec:biblio_microstructure});
    central case study throughout Chs~23--26.
  \item[\citet{prosperity4wiki} --- IMC Prosperity 4 wiki.] The
    primary source of truth for the current competition's rules,
    mechanics, and round schedule. Cited in Chs~23--26.
  \item[\citet{jmerle_backtester} --- Merle's Prosperity 3
    backtester.] An open-source local backtester for Prosperity
    algorithms; indispensable for iterating between rounds. Cited
    in Chs~23, 24.
  \item[\citet{jmerle_visualiser} --- Merle's Prosperity 3
    visualiser.] The companion visualiser; lets you debug why your
    algorithm got filled where it did. Cited in Ch~23.
  \item[\citet{nash1950bargaining} --- Nash bargaining.] The game
    theory paper behind the Prosperity ``Shells from Sea Turtles''
    manual-trading round. Cited in Ch~26.
  \item[\citet{thaler1988anomalies} --- Thaler on the winner's curse.]
    Explains why sealed-bid auctions systematically overpay; relevant
    to Prosperity's manual rounds. Cited in Ch~26.
\end{description}

\subsection{Rates and fixed income}

\begin{description}
  \item[\citet{hullwhite1990} --- Hull and White, 1990.] The
    short-rate model that is still the front-office workhorse for
    Bermudan swaption pricing. Cited in Chs~28, 29.
  \item[\citet{hagan2002managing} --- Hagan, Kumar, Lesniewski,
    Woodward, \emph{Managing Smile Risk}.] The foundational SABR
    paper; every rates vol trader knows the asymptotic expansion by
    heart. Cited in Chs~17, 29.
  \item[\citet{braceetal1997} --- Brace, G\k{a}tarek and Musiela.]
    The LIBOR Market Model (BGM/LMM) paper; how to price exotics
    consistent with the full swaption surface. Cited in Ch~29.
  \item[\citet{hullwhite2012fva} --- Hull and White, FVA debate.]
    The original ``funding is not a valuation adjustment'' argument;
    you must know both sides for any Strats interview. Cited in
    Chs~33, 37.
\end{description}

\subsection{Credit}

\begin{description}
  \item[\citet{li2000default} --- Li's Gaussian copula.] The paper
    that industrialised CDO pricing and, two decades later, was
    blamed for the financial crisis. Cited in \cref{ch:credit_derivatives}
    (Ch~29/30).
  \item[\citet{salmon2009formula} --- Salmon, Wired article.] The
    ``Formula That Killed Wall Street'' --- an accessible and
    deliberately provocative retelling of the Li copula story.
    Cited in Ch~29/30.
\end{description}

\subsection{FX and commodities}

\begin{description}
  \item[\citet{garman1983foreign} --- Garman and Kohlhagen.] The
    FX version of Black--Scholes (two rates instead of one). Cited
    in Ch~30.
  \item[\citet{clark2011foreign} --- Clark, \emph{FX Option Pricing}.]
    The practitioner's FX vol reference; covers premium-adjusted
    deltas, smile conventions, and barrier pricing. Cited in Ch~30.
  \item[\citet{castagna2007consistent} --- Castagna, \emph{FX Options
    and Smile Risk}.] The Vanna--Volga construction for FX smile;
    the market standard for interpolation. Cited in Ch~30.
  \item[\citet{borio2016covered} --- BIS, Covered Interest Parity
    Lost.] Documents the post-2008 breakdown of CIP; essential
    modern FX reading. Cited in Chs~13, 30.
  \item[\citet{schwartz1997stochastic} --- Schwartz 1997.] The
    three-factor commodity pricing paper; where mean reversion in
    convenience yields was established. Cited in Ch~31.
  \item[\citet{gibsonschwartz1990} --- Gibson and Schwartz.] The
    earlier two-factor oil model; foundation of Ch~31's commodity
    curve material.
  \item[\citet{schwartzsmith2000} --- Schwartz and Smith.] The
    short-term/long-term decomposition used on commodity desks
    today. Cited in Ch~31.
\end{description}

\subsection{Numerical methods}

\begin{description}
  \item[\citet{glasserman2004monte} --- Glasserman, \emph{Monte
    Carlo Methods in Financial Engineering}.] The standard MC
    reference; variance reduction, LSM for Americans, sensitivities.
    Cited in Chs~34, 17.
  \item[\citet{longstaffschwartz2001} --- Longstaff and Schwartz.]
    The LSM regression method for pricing American options by
    simulation; the standard approach on every rates and equities
    exotics desk. Cited in Ch~34.
  \item[\citet{cranknicolson1947} --- Crank and Nicolson.] The 1947
    paper that gave the finite-difference scheme used for pricing
    any PDE-based option model. Cited in Chs~34, App~A.
  \item[\citet{coxrossrubinstein1979} --- CRR 1979.] See options
    section; also a cornerstone numerical method.
\end{description}

\subsection{Interviews}

\begin{description}
  \item[\citet{zhou2020quant} --- Zhou, the ``Green Book''.] Hundreds
    of interview problems with full solutions; the book every quant
    candidate has actually read, and the source of \cref{ch:interview_problems}'s
    problem selection. Cited in Ch~38.
\end{description}

\subsection{GS culture, institutional context, and careers}

\begin{description}
  \item[\citet{derman2004mylife} --- Derman, \emph{My Life as a
    Quant}.] Emanuel Derman's memoir of physics-to-quant; the
    flavour of a 1990s trading floor, and where the BDT and local-vol
    models came from. Cited in Chs~17, 27, 36.
\end{description}

\subsection{Mathematical background}

\begin{description}
  \item[\citet{grimmett2001probability} --- Grimmett and Stirzaker.]
    The standard UK undergraduate probability book; measure-theoretic
    from Ch 5 onwards, with excellent exercises. Cited in
    \cref{ch:app_math} (App~A).
  \item[\citet{strang2016linalg} --- Strang, linear algebra.]
    The best introductory linear-algebra text for intuition; pairs
    with MIT OCW 18.06 lectures. Cited in App~A.
  \item[\citet{trefethen1997numerical} --- Trefethen and Bau.] The
    gold-standard book for numerical linear algebra; QR, SVD, and
    conditioning done properly. Cited in App~A, Ch~34.
  \item[\citet{oksendal2003stochastic} --- {\O}ksendal.] The
    standard first book on SDEs and It\^o calculus; readable,
    complete proofs, exercises. Cited in App~A, Chs~14, 16, 34.
  \item[\citet{shreve2004stochastic} --- Shreve Vol.~II.] The
    de facto Masters-in-Financial-Mathematics textbook; rigorous
    continuous-time mathematical finance. Cited in App~A, Chs~14--17,
    34.
  \item[\citet{pham2009continuous} --- Pham.] Stochastic control
    and HJB equations applied to finance; the heavy machinery behind
    Avellaneda--Stoikov and Almgren--Chriss. Cited in Chs~8, 18, 34.
  \item[\citet{duffy2006finite} --- Duffy.] The practical guide to
    finite-difference methods for derivatives pricing; Crank--Nicolson,
    operator splitting, American options. Cited in Ch~34.
\end{description}

% =========================================================================
\section{Primary sources that are free online}
\label{sec:app_biblio_free}
% =========================================================================

A non-trivial fraction of the best quant-trading material is free.
This section collects the items from the bibliography above that can
be read legally without paying for a textbook or a journal
subscription. For a self-funded reader this is the shortest path to
covering most of the book's territory.

\paragraph{MIT OpenCourseWare --- 18.S096 \emph{Topics in Mathematics
with Applications in Finance}.} Free video lectures and PDF notes from
2013, covering measure theory, It\^o calculus, Black--Scholes,
stochastic volatility, and portfolio theory. Lectures~15--21 map
almost one-to-one onto \cref{ch:black_scholes}--\cref{ch:vol_surface}
of this book. Available at \url{https://ocw.mit.edu/courses/18-s096-topics-in-mathematics-with-applications-in-finance-fall-2013/}.

\paragraph{Bennett, \emph{Trading Volatility}.} A 250-page sell-side
research document distributed free as a PDF at
\url{https://www.trading-volatility.com/}. The single best free
reference for equity vol trading. A long-standing mirror also exists
via Google Scholar searches for the title.

\paragraph{Rosenthal, UIC lecture slides on market microstructure.}
Hosted on the University of Illinois Chicago course pages for the
Financial Engineering master's programme; each lecture is a
self-contained PDF on one microstructure topic (matching engines,
spread estimation, market impact).

\paragraph{Damodaran's option arbitrage and investment fables notes.}
NYU Stern hosts Aswath Damodaran's entire library of teaching
materials at \url{https://pages.stern.nyu.edu/~adamodar/}. The
options arbitrage note (\citealp{damodaranOptionsArbitrage}) is a
particularly clean, free summary of put--call parity and related
bounds.

\paragraph{QuantStart, Hummingbot, Letian Wang, and Reasonable
Deviations blogs.} Practitioner-level tutorials with working code.
QuantStart covers backtesting and HMM regime detection
(\citealp{quantstart2016hmm}); Hummingbot publishes the best free
walk-through of Avellaneda--Stoikov (\citealp{hummingbotASguide});
Letian Wang's blog has a long series on pairs trading and cointegration;
Reasonable Deviations summarises L\'opez de Prado's \emph{AFML}
chapter by chapter. The URLs are in the \texttt{.bib} file.

\paragraph{arXiv and SSRN.} Most of the cited journal papers
(\citealp{avellaneda2008high}, \citealp{almgren2000optimal},
\citealp{avellaneda2010stat}, \citealp{bailey2014deflated},
\citealp{moskowitz2012time}) are on arXiv or SSRN as author
preprints. Search for the paper title; the author's institutional
page is usually the first hit and legally hosts a free copy.

\paragraph{Pelgrims, \emph{Matching Engines}.} A clear blog-post
walk-through of a price-time-priority matching engine's internals
(\citealp{pelgrims2020matching}) at
\url{https://jellepelgrims.com/posts/matching_engines}; the best
free companion to \cref{ch:lob}.

\paragraph{Prosperity 4 wiki.} The official competition wiki
(\citealp{prosperity4wiki}) is free and is the canonical source for
rules, mechanics, and the tutorial-round simulator. The Merle
backtester and visualiser (\citealp{jmerle_backtester};
\citealp{jmerle_visualiser}) are open-source on GitHub.

\paragraph{MIT OpenCourseWare --- 18.06 Linear Algebra (Strang).}
Free companion lectures to \citet{strang2016linalg}; the best way to
rebuild linear-algebra fluency from scratch before Chapters 11, 17,
and 34.

% =========================================================================
\section{What's not in this book}
\label{sec:app_biblio_gaps}
% =========================================================================

No single book can cover every part of quantitative finance. This one
deliberately stops where depth would cost the reader more than it
gave back, or where the material is so specific to a single firm or a
single piece of hardware that a textbook would be the wrong vehicle
for it. Here are the conscious omissions, with a pointer to where
you would go next if you need them.

\begin{description}
  \item[Microsecond-latency execution engineering.] This book covers
    execution cost models (\citealp{almgren2000optimal},
    \citealp{almgren2005direct}) and order-book dynamics, but it does
    not teach you how to build an exchange gateway that quotes at
    500~ns. That territory is FPGA design, kernel bypass (Solarflare
    OpenOnload, DPDK), and colocation-rack physics. Start with
    Donald MacKenzie's sociology-of-HFT papers, the Optiver/Jane
    Street engineering blogs, and the talks on the CppCon YouTube
    channel under ``low latency''. Nothing on the GS Off-Cycle
    placement requires this.
  \item[Regulatory compliance at operational depth.] Chapter~37
    covers the capital-regulation frameworks (Basel III,
    Volcker, Dodd--Frank, MiFID~II) at the level a Strat needs to
    understand a trading restriction. It does not teach you how
    to \emph{be} the compliance officer who implements them. For
    that, the Financial Conduct Authority's FCA Handbook and the
    Office of the Comptroller of the Currency's guidance are the
    primary sources. Read ICAAP and ILAAP documentation from any
    large UK bank's pillar~3 disclosures for practical flavour.
  \item[Case-by-case firm history beyond LTCM, MG, and Amaranth.]
    The blow-ups covered in Chapter~19/35 are LTCM
    (\citealp{lowenstein2000ltcm}), Metallgesellschaft, and
    Amaranth, because they illustrate different failure modes
    (leverage/correlation, hedging-rollover, concentration).
    Knight Capital 2012, Archegos 2021, and the 2020
    WTI-goes-negative episode are all worth reading up on. Gregory
    Zuckerman's \emph{The Man Who Solved the Market} (on Renaissance)
    and Patterson's \emph{The Quants} fill in the firm-by-firm
    picture.
  \item[Crypto-specific microstructure.] Centralised-exchange crypto
    (Binance, Coinbase, FTX-before-2022) behaves \emph{almost} like
    traditional equity markets, with tick sizes and limit order
    books. Decentralised-exchange (AMM) microstructure --- Uniswap
    v2/v3, constant-product curves, impermanent loss, MEV, sandwich
    attacks --- is genuinely different and is not covered here. The
    Paradigm Research and a16z crypto research blogs are the right
    next step.
  \item[Specific firm stacks beyond SecDB/Slang.] Chapter~36
    deliberately drills into SecDB and Slang because they are
    Goldman-specific and directly relevant to the Off-Cycle role.
    The equivalent stacks at JP~Morgan (Athena/Python), Morgan
    Stanley (MSS-inhouse), Bank of America (Quartz/Python), Deutsche
    Bank (gRisk/Java), and Citadel (C++/Python) get at most a
    sentence. If you move firm, your first two weeks will replace
    this chapter with an internal wiki.
  \item[Tail-risk and catastrophe modelling at insurance depth.]
    VaR and expected shortfall (Chapter~19/35) are as far as this
    book goes. Full tail-risk modelling --- extreme value theory,
    peaks-over-threshold, copulas for dependence in the tail --- is
    covered in McNeil, Frey and Embrechts, \emph{Quantitative Risk
    Management}, which is the standard next step.
  \item[Behavioural finance beyond a single Thaler reference.] The
    winner's curse (\citealp{thaler1988anomalies}) gets a citation;
    prospect theory, herding, disposition effect, and the full
    Kahneman--Tversky programme do not. Kahneman's \emph{Thinking,
    Fast and Slow} is the accessible starting point.
\end{description}

\begin{pitfallbox}[Don't treat this book as comprehensive]
This book is a \emph{starting point} for two specific goals: the IMC
Prosperity competition and a Goldman Sachs Off-Cycle placement. It is
not an encyclopaedia of quantitative finance and it does not try to
be. Two failure modes to avoid:
\begin{enumerate}
  \item \emph{Assuming that what's in here is all there is.} If you
    reach a production desk and find yourself saying ``the textbook
    didn't mention FVA/XVA/cross-currency basis/repo squeezes'', that
    is the textbook working as intended. The next step is the
    pillar~3 disclosures, Hull's Chapter~33 on XVAs, and your desk's
    internal documentation --- none of which would fit in a single
    book.
  \item \emph{Assuming that a cited result is still current.} Market
    microstructure papers from the 1980s (\citealp{kyle1985continuous},
    \citealp{glosten1985bid}) remain conceptually correct but
    describe a world of floor-traded specialists; their empirical
    magnitudes (spread sizes, impact coefficients) are out of date
    by orders of magnitude. Always cross-check with
    \citet{bouchaud2018trades} or the latest BIS Quarterly Review
    for modern numbers.
\end{enumerate}
When in doubt, follow the citation trail. Every topic in this book
has at least one primary source listed in the bibliography above; the
primary source will always be deeper, and usually more current, than
the summary offered here. The point of this book is to make that
primary source readable --- not to replace it.
\end{pitfallbox}

That is the end of the book. If you have reached this page having
worked through one of the three reading orders above, you are now at
the level where the bibliography, not the chapters, is the next step.
Good luck in Prosperity, and good luck on the desk.
